%% kapa: exam=1 round=32 year=2021 period=1 form=A subject=real_estate subject_ko=부동산학원론 seq=34 number=114
%% source: Q-Net 감정평가사 기출문제 (cst00309) artlSeq=5208397 file=extracted/1008/5208397/A형/2021년도 제32회 감정평가사 1차 1교시 A형.pdf
%% conversion: kordoc 4.24.51 -> tools/gen1.py (auto); review=proofread-sonnet; flags=none
\begin{question}{114}
감정평가 과정상 지역분석과 개별분석에 관한 설명으로 옳지 않은 것은?
\begin{choices}
  \item 지역분석을 통해 해당 지역 내 부동산의 표준적 이용과 가격수준을 파악할 수 있다.
  \item 지역분석은 개별분석보다 먼저 실시하는 것이 일반적이다.
  \item 인근지역이란 대상부동산이 속한 지역으로서 부동산의 이용이 동질적이고 가치형성요인 중 개별요인을 공유하는 지역을 말한다.
  \item 유사지역이란 대상부동산이 속하지 아니하는 지역으로서 인근지역과 유사한 특성을 갖는 지역을 말한다.
  \item 지역분석은 대상지역에 대한 거시적인 분석인 반면, 개별분석은 대상부동산에 대한 미시적인 분석이다.
\end{choices}
\end{question}
\answer{3}
